\subsection{Complex heat time and real-time translations}
\label{sec:complex-heat}
The complex heat extension used in the boundary argument can be constructed
from the real Euclidean semigroup alone. This also fixes the sign of the
real-time translation group. Let $(E_t)_{t\geq0}$ be a strongly continuous
semigroup of self-adjoint contractions on a complex Hilbert space, with
$E_0=I$. Each $E_t$ is positive because $E_t=E_{t/2}^2$.

We use continuous functional calculus for a bounded self-adjoint operator:
continuous scalar functions on its spectrum act as a unital isometric
star homomorphism, with the usual composition rule. This is the
self-adjoint instance of the continuous calculus in \cite[Theorem 2.19]{Williams}.
We use the continuous calculus for positive roots and the bounded Borel calculus for the imaginary powers.

\begin{proposition}\label{prop:complex-heat}
There is a holomorphic operator-valued function $E(z)$ on
$\operatorname{Re}z>0$ such that
\[
 \|E(z)\|\leq1,\quad E(z+w)=E(z)E(w),\quad
 E(z)^*=E(\overline z),\quad E(t)=E_t\quad(t>0).
\]
Each $E_t$, $t>0$, is injective and has dense range. There is a strongly
continuous unitary group $V(t)$, $t\in\mathbb R$, characterized by
\begin{equation}\label{eq:complex-heat-edge}
 V(t)E(z)=E(z)V(t)=E(z-it)\qquad(\operatorname{Re}z>0).
\end{equation}
Every bounded operator commuting with all $E_t$ commutes with $E(z)$ and
$V(t)$. If $E_t\Omega=\Omega$ for all $t$, then
$E(z)\Omega=V(t)\Omega=\Omega$.
\end{proposition}
\begin{proof}
If $E_tx=0$ for some $t>0$, self-adjointness and the semigroup law
give $E_{t/2}x=0$. Iterate and use strong continuity at zero to get
$x=0$. Since $E_t$ is self-adjoint, its range is dense.

Write $A=E_1$. By the bounded Borel functional calculus
\cite[Theorem~VII.2]{ReedSimonI} the spectral measure of $A$ is
supported on $[0,1]$ and has zero projection at $0$ by injectivity.
Define $E(z)=A^z$ and $V(t)=A^{-it}$ using the real logarithm on
$(0,1]$. The values of the multipliers at zero are immaterial.
The product, adjoint, norm and edge identities follow from the
scalar multipliers. Dominated convergence for the vector spectral
measures makes $V(t)$ strongly continuous. On each compact
subset of $\operatorname{Re}z>0$, the bound on
$\lambda^\delta|\log\lambda|^2$ for any $\delta>0$ gives the
uniform Taylor remainder needed for norm holomorphy.

For positive integers $n$, $E_{1/n}$ is the positive $n$th root of
$A$ by the semigroup law and the continuous calculus. Thus
$E(p/n)=E_{p/n}$ for positive rational $p/n$. Norm continuity of
$E(z)$ and strong continuity of the semigroup identify them on
all positive real times. A bounded operator commuting with the
semigroup commutes with $A$ and hence its Borel calculus. If
$A\Omega=\Omega$, the scalar multipliers at $1$ give both vacuum
identities. Finally, the edge identity on the dense range of $A$
determines $V(t)$.
\end{proof}

\subparagraph{Coordinate convention in Sections~\ref{sec:complex-heat}--\ref{sec:locality}.}
Section~\ref{ex-statement} writes points of $\mathbb R^4$ as
$x=(\boldsymbol x,t)$, with time last. In
Sections~\ref{sec:complex-heat}--\ref{sec:locality} we write instead
$a=(a^0,\mathbf a)$ and $X_i=(t_i,x_i)$ with time first, so that index $0$
is the time coordinate, and matrices such as
$\eta=\operatorname{diag}(1,-1,-1,-1)$ are written in this order. The two
descriptions are related by the linear change of coordinates
$C(t,\mathbf x)=(\mathbf x,t)$ from time-first to time-last coordinates,
with $|\det C|=1$ and $C^{-1}=C^{\mathsf T}$. This is a change of
realization, not a symmetry of either one. The Minkowski matrices satisfy
$\eta_{\mathrm{last}}=C\eta_{\mathrm{first}}C^{\mathsf T}$; a Lorentz matrix
$\Lambda$ in time-first coordinates corresponds to $C\Lambda C^{-1}$; the
time-last translation by $a$ corresponds to the time-first translation by
$C^{-1}a$; and the forward cone $\{p^0\geq|\mathbf p|\}$ corresponds to
$\{(\mathbf p,E):|\mathbf p|\leq E\}$. Accordingly, the statements of
these sections transfer to the time-last convention of
Section~\ref{ex-statement} under the following substitutions: test
functions are pulled back, $f\mapsto f\circ C$; the Poincar\'e action is
conjugated by $C$; and supports, cones and spectral measures are pushed
forward by $C$. The pullback is a topological linear isomorphism of
$\mathcal S(\mathbb R^4)$ commuting with complex conjugation. It preserves
spacelike separation of supports for the corresponding forms, and integrals
acquire no Jacobian factor.

Suppose now that spatial translations $U(x)$ form a strongly continuous
unitary representation commuting with the heat semigroup and fixing
$\Omega_E$. The proposition
defines a strongly continuous unitary representation of all real translations,
\[
 T_E(a^0,\mathbf a)=V(a^0)U(\mathbf a).
\]
The group law follows from commutation. Joint strong continuity follows by
adding and subtracting $V(a^0)U(\mathbf b)\xi$ and using the unitary norms.
Moreover $E(\varepsilon-it)\xi=V(t)E_\varepsilon\xi\to V(t)\xi$ as
$\varepsilon\downarrow0$, which justifies the boundary notation $E(-it)$
used in the heat-vector chapter. No operator-norm continuity at the boundary
is required.

\begin{lemma}\label{lem:time-fixed-heat}
The vectors fixed by every $V(t)$ are exactly the vectors fixed by every
$E_s$. Consequently, if the selected heat ground space is
$\mathbb C\Omega_E$, the same is true of the subspace fixed by all real
translations.
\end{lemma}
\begin{proof}
A heat-fixed vector satisfies $A\xi=\xi$, so functional calculus and the
definition of $V$ give $V(t)\xi=\xi$. Conversely suppose $V(t)\xi=\xi$.
The holomorphic vector function $g(z)=E(z)\xi$ satisfies
$g(z-it)=E(z)V(t)\xi=g(z)$. Its derivative in the imaginary direction is
zero; complex differentiability therefore makes its full derivative zero.
It is constant on the connected right half-plane. Since
$E_s\xi\to\xi$ as $s\downarrow0$, this constant is $\xi$.
A vector invariant under all translations is, in particular, invariant under
$V(t)$, and hence belongs to the heat ground space. The vacuum itself is
translation-invariant, which proves the last assertion.
\end{proof}

Combining this lemma with the isometric translation intertwiner in
Proposition~\ref{prop:word-motion} proves uniqueness of the invariant vacuum
in the reconstructed field Hilbert space: an invariant $\xi$ maps to
$c\Omega_E=\mathcal V(c\Omega)$, and injectivity of the intertwiner
$\mathcal V$ gives $\xi=c\Omega$.
